\section*{Practical Appendices: Financial Instruments for Practice}

\subsection*{Appendix A: The 13-Week Cash Forecast}

Use a spreadsheet with 1 column per week and a final column for the person responsible. A forecast is useful when the founder can reconcile it to the bank and identify the \textbf{1st date of an unacceptable cash balance}. If receipts are uncertain, show at least base and delayed-collection scenarios. This worksheet is a managerial forecast rather than a statutory cash-flow statement.

\begin{table*}[htbp]
\centering
\small
\begin{tabularx}{\textwidth}{Xrrrrp{3.5cm}}
\toprule
\textbf{Forecast Line (Illustrative Category)} & \textbf{Week 1} & \textbf{Week 2} & \textbf{\dots} & \textbf{Week 13} & \textbf{Owner / Source} \\
\midrule
Opening reconciled available cash & & & & & Bank statement \\
Collected customer cash: contracted & & & & & Contract / invoice register \\
Collected customer cash: speculative & & & & & Pipeline (risk-discounted) \\
Refunds and chargebacks & & & & & Policy / student claims \\
Personnel and contractor payments & & & & & Payroll / contractor terms \\
Platform, rent, and infrastructure & & & & & Approved commitments \\
Tax, statutory, and insurance & & & & & Compliance calendar \\
Loan interest and principal & & & & & Lender amortization table \\
Equipment and capital spending & & & & & Approved project budget \\
Other actual disbursements & & & & & Source invoices \\
\midrule
\textbf{Closing cash} & & & & & Reconciled formula \\
\textbf{Minimum acceptable cash} & & & & & Reserve floor policy \\
\midrule
\textbf{Headroom / Shortfall} & & & & & Action trigger \\
\bottomrule
\end{tabularx}
\caption{The 13-week rolling cash forecast architecture identifying liquidity troughs.}
\end{table*}

\textbf{Weekly operating rule:} Replace predictions with actual transactions and roll the forecast forward 1 week. Never move an overdue receipt to a later column without preserving its previous forecast date and the reason for slippage. Keep tax and refund forecasts separately identifiable. Treat projected new loans or equity as contingent until contractually available. Where cash is restricted, exclude it from \emph{available} cash rather than counting it as operational headroom.

The most useful output is often a short narrative: ``Our minimum forecast balance is ₹X in week Y. It falls below our reserve policy if Client Z delays payment by 2 weeks. Our preventive action is to obtain invoice acceptance now and defer a discretionary vendor commitment.'' The actual numbers must come from the venture's records.

\subsection*{Appendix B: The 3-Statement Reconciliation Checklist}

Before claiming the books are ready, verify the following 10 relationships:
\begin{enumerate}
\item \textbf{Closing balance sheet identity:} total assets equal total liabilities plus equity, including retained results and documented capital changes.
\item \textbf{Cash bridge:} opening cash plus net cash movements equals closing cash, with relevant bank reconciliation adjustments properly recorded.
\item \textbf{Revenue bridge:} amounts received before performance are treated under applicable rules as contract liabilities or another appropriate balance, rather than being added to earned revenue automatically.
\item \textbf{Receivables bridge:} opening receivables plus amounts billed or earned as applicable minus collections and adjustments explain closing receivables.
\item \textbf{Payables bridge:} costs incurred but unpaid appear in an appropriate payable or accrual rather than disappearing from profit.
\item \textbf{Capital expenditure bridge:} cash used for qualifying equipment appears as asset purchases, while depreciation or amortization is recognized over the relevant periods.
\item \textbf{Debt bridge:} loan proceeds and principal repayment change the borrowing balance; interest and financing fees are recorded under applicable rules.
\item \textbf{Equity bridge:} capital contributions and distributions are distinguished from revenues and expenses.
\item \textbf{Operating-cash reconciliation:} appropriate profit measure, noncash items, and changes in operating balances explain the operating cash-flow figure under the chosen standard.
\item \textbf{Contracts and promises:} the amount and expected cash cost of undelivered services are reviewed alongside the accounting contract balance: these are related but not identical measurements.
\end{enumerate}

\textbf{Worked checkpoint:} In the main fictional 1st-quarter case:
\begin{itemize}
\item Closing cash of ₹8,99,000 plus receivables ₹60,000 plus net equipment ₹1,00,000 equals assets of \textbf{₹10,59,000}.
\item Liabilities ₹4,65,000 (loan ₹2,00,000 $+$ reviewer payable ₹25,000 $+$ contract liability ₹2,40,000) plus equity ₹5,94,000 (contributed ₹6,00,000 less loss ₹6,000) agree at \textbf{₹10,59,000}.
\item Operating cash of ₹2,19,000 less equipment ₹1,20,000 plus financing ₹8,00,000 equals the \textbf{₹8,99,000} closing cash figure.
\end{itemize}
The entire financial story reconciles.

\subsection*{Appendix C: A 1-Page Working-Capital Dossier}

For each material commercial route fill in:

\begin{quote}
\small
\textbf{OFFER / CUSTOMER TYPE:}\\
\textbf{STANDARD PRICE AND ACTUAL COLLECTED PRICE:}\\
\textbf{CONTRACT SIGNATURE AND PURCHASE AUTHORIZATION:}\\
\textbf{INVOICING MILESTONES:}\\
\textbf{COLLECTION DATES AND ACTUAL AGING:}\\
\textbf{INITIAL CASH NEEDED BEFORE 1ST COLLECTION:}\\
\textbf{REMAINING CASH REQUIRED TO FULFILL THE PROMISE:}\\
\textbf{EXPECTED REFUNDS / CREDITS / WITHHOLDINGS:}\\
\textbf{PAYABLES BY DUE DATE:}\\
\textbf{INVENTORY OR EQUIPMENT INVESTMENT, IF ANY:}\\
\textbf{LIKELY CASH TROUGH AND DATE:}\\
\textbf{MAXIMUM CAPITAL TIED UP UNDER BASE CASE:}\\
\textbf{MAXIMUM CAPITAL TIED UP UNDER DELAYED CASE:}\\
\textbf{ALTERNATIVE TERMS TO NEGOTIATE:}\\
\textbf{DECISION OWNER AND NEXT ACTION:}
\end{quote}

For a college or enterprise contract, obtain documentary confirmation of the contracting entity, authorized budget, acceptance process, and invoice recipient. For a cohort, distinguish refundable deposits, tuition collected, tuition earned, and the cash cost of undelivered sessions. For physical goods, track inventory purchases and supplier payments separately from customer receipts. The aim is to identify the financing need created by \emph{1 additional sale}.

\subsection*{Appendix D: The Break-Even and Capacity-Band Worksheet}

Use the following calculations only within their stated assumptions:
\begin{itemize}
\item Realized revenue per customer: $P$.
\item Incremental cost per customer, including honestly attributable acquisition and support: $v$.
\item Fixed costs within the relevant capacity band: $F$.
\item Simple unit contribution: $P - v$.
\item Simple break-even units, if $P > v$: $\lceil F / (P - v) \rceil$.
\item Result for $Q$ units within the same band: $Q(P - v) - F$.
\end{itemize}

Now add the operational reality:

\begin{table}[htbp]
\centering
\small
\begin{tabularx}{\columnwidth}{p{2.2cm}Xp{2.0cm}}
\toprule
\textbf{Capacity Band} & \textbf{Workload Limit \& Fixed Step} & \textbf{Pre-Receipt Cash Needed?} \\
\midrule
Current resources & Up to 16 learners; $F = ₹2,20,000$ & Minimal \\
1st expansion step & 17 to 30 learners; $+₹60,000$ reviewer & 1 month mentor advance \\
2nd expansion step & 31 to 50 learners; $+₹1,20,000$ senior lead & 2 months faculty salary \\
\bottomrule
\end{tabularx}
\caption{Capacity bands mapping fixed step-costs and pre-receipt cash requirements.}
\end{table}

\textbf{Illustrative trap:} With $P = ₹30,000$, $v = ₹8,000$, $F = ₹2,20,000$, and maximum capacity 16, operating result at 16 is ₹1,32,000. If the 17th participant triggers ₹60,000 of additional fixed cost, result at 17 becomes ₹94,000. The arithmetic is not telling you to reject the 17th learner forever: it tells you to design the staffing step and enrollment plan rather than promise a smooth cost curve.

Repeat the analysis after assigning realistic replacement costs to founder teaching, sales, and administration. If the normalized economics fail, state that explicitly even when actual cash expenditure remains low.
